% Measuring What Kills Your App: A Methodology for Detecting Background
% Execution Failures in Production Android Apps
% Target: IEEE Software (rolling submissions, ~4000-6000 words)
% Format: IEEEtran journal/compsoc, single-anonymous

\documentclass[10pt,journal,compsoc]{IEEEtran}

\usepackage{booktabs}
\usepackage{listings}
\usepackage{xcolor}
\usepackage{url}
\usepackage{cite}
\usepackage{enumitem}

\lstset{
  basicstyle=\ttfamily\small,
  breaklines=true,
  frame=single,
  language=Java,
  commentstyle=\color{gray},
  keywordstyle=\color{blue},
  numbers=left,
  numberstyle=\tiny\color{gray},
  xleftmargin=2em,
}

\begin{document}

\title{Measuring What Kills Your App:\\A Methodology for Detecting Background Execution Failures\\in Production Android Apps}

\author{Sagar~Ghag
\IEEEcompsocitemizethanks{\IEEEcompsocthanksitem S. Ghag is with Kryonex Labs, Pune, India.\protect\\
E-mail: sagar.ghag1997@gmail.com}}

\IEEEtitleabstractindextext{%

% ============================================================================
% ABSTRACT
% ============================================================================

\begin{abstract}
Background location tracking apps die in the field. Doze mode suspends
them. OEM battery killers terminate their foreground services. Users
revoke location permissions mid-shift. Practitioners know this from
dontkillmyapp.com and production firefighting, but there is no published
methodology for systematically detecting, classifying, and measuring
these failures in a running app. The result: teams build ad-hoc
telemetry that catches some failures and misses others, with no way to
know which.

We present a measurement methodology for background execution failures
on Android. The methodology comprises three parts: a kill-cause taxonomy
that classifies failures into five categories with distinct observable
signals, an instrumentation specification that defines the minimum
telemetry a production app needs to detect each category, and a
classification procedure that maps observed signals to kill causes. We
validate the methodology against a production Flutter app (104,886 lines,
86 physical devices across three organisations (two accounting for 93\% of the
fleet), two months of tracked field
deployment) by auditing its existing telemetry against the specification.
The audit found six detection signals already present and twelve gaps
where failures go unrecorded. A database analysis of 159 tracking
sessions and 14,094 waypoints revealed 509 timestamp gaps exceeding five
minutes, 65 sessions (40.9\%) that captured three or fewer waypoints
across multi-hour sessions, and 25 sessions (15.7\%) that ran until
server timeout with no user-initiated end. The most consequential
instrumentation gap: process kills that occur during a permission
revocation are invisible because the revocation state is held in memory
and lost with the process.

The methodology is framework-agnostic in its taxonomy and classification
rules. The instrumentation examples use Flutter and Android, but the
signals and failure modes apply to any app that runs background services
on Android.
\end{abstract}

\begin{IEEEkeywords}
Android, background execution, process kills, permission revocation,
mobile telemetry, measurement methodology, OEM fragmentation.
\end{IEEEkeywords}}

\maketitle

% ============================================================================
% 1. INTRODUCTION
% ============================================================================

\section{Introduction}
\label{sec:intro}

\IEEEPARstart{A}{~tracking} app that dies in the background looks the same
as a tracking app that was never running. No crash report. No ANR dialog. The
foreground service stops, the GPS radio goes silent, and three hours
later the user opens the app to find a gap in the data where a
delivery route should be.

This is the central measurement problem for any Android app that depends
on sustained background execution: location trackers, health monitors,
fleet management tools, fitness apps. The OS actively works to kill them.
Stock Android has Doze mode~\cite{android_doze} and App Standby
Buckets~\cite{android_buckets}. OEM forks go further: Huawei's
PowerGenie, Xiaomi's MIUI Autostart restrictions, OnePlus Deep
Optimization, and Samsung's Sleeping Apps each enforce their own kill
policies independently of stock Android
behaviour~\cite{dontkillmyapp}. Users add a third vector by revoking
permissions, sometimes intentionally, sometimes because Android 12+
auto-revokes permissions for apps it considers
inactive~\cite{android_hibernate}.

The practitioner community tracks these problems through
dontkillmyapp.com, a crowdsourced catalogue of OEM kill behaviours. The
site is useful but unstructured: it lists workarounds per OEM without
a detection methodology, without severity quantification, and without a
way for an app to measure its own kill rate in the field.

Academic work on mobile background execution focuses on energy
consumption~\cite{li2014making,pathak2012energy} and permission
models~\cite{felt2012android,au2012pscout}. The gap is in the middle:
how to \emph{observe} what kills a background app in production, at what
rate, and by what cause. This is the gap we address.

We contribute:

\begin{enumerate}
  \item A \textbf{kill-cause taxonomy} that classifies background
        execution failures into five categories, each with distinct
        observable signals that a production app can detect
        (Section~\ref{sec:taxonomy}).
  \item An \textbf{instrumentation specification} that defines the
        minimum set of events, metadata, and persistence requirements
        for detecting each failure category
        (Section~\ref{sec:instrumentation}).
  \item A \textbf{classification procedure} that maps observed telemetry
        signals to kill causes, including rules for disambiguating
        overlapping signals (Section~\ref{sec:classification}).
\end{enumerate}

We validate the methodology against a production Flutter app with 86
physical devices across three organisations, auditing its existing ad-hoc
telemetry against the specification (Section~\ref{sec:validation}).

% ============================================================================
% 2. BACKGROUND: THE KILL LANDSCAPE
% ============================================================================

\section{How Android Kills Background Apps}
\label{sec:landscape}

Android's relationship with background apps has grown adversarial since
API 23. Each OS version adds new restrictions; each OEM fork adds
vendor-specific ones on top. For a developer building a location tracking
app, ``Android'' is not one platform. It is at least six, and three of
them will kill your foreground service without telling you.

\subsection{Stock Android Restrictions}

Doze mode (API 23+) suspends network access and defers alarms when the
device is stationary and screen-off~\cite{android_doze}. Light Doze
(API 24+) activates sooner, within minutes of screen-off on some
devices. App Standby Buckets (API 28+) throttle background work based on
recency of use: an app in the ``rare'' bucket gets fewer than 5 minutes
of background execution per day~\cite{android_buckets}. Background
Location Limits (API 26+) reduce location update frequency to a few
times per hour for apps without a foreground service.

These are documented. An app can request
\texttt{REQUEST\_IGNORE\_BATTERY\_OPTIMIZATIONS} and run a foreground
service with \texttt{FOREGROUND\_SERVICE\_LOCATION} (API 34+).  In
theory, this exempts it from all of the above. In practice, it exempts
it from stock Android's restrictions. OEM restrictions are a separate
layer.

\subsection{OEM Kill Policies}

Table~\ref{tab:oem} summarises the OEM kill behaviours we have observed
in production across six months of deployment on budget and mid-range
Android devices in an emerging market where budget handsets dominate.

\begin{table}[h]
\caption{OEM background kill behaviours observed in production}
\label{tab:oem}
\small
\begin{tabular}{@{}p{1.6cm}cp{4.8cm}@{}}
\toprule
OEM & Aggression & Observed behaviour \\
\midrule
Huawei    & 5/5 & PowerGenie kills foreground services independently of battery optimisation whitelist. Survives \texttt{isIgnoring\-BatteryOptimizations} check. \\
Xiaomi    & 5/5 & MIUI Autostart disabled by default. Service killed within minutes of screen-off unless manually whitelisted. Settings reset on MIUI updates. \\
OnePlus   & 5/5 & Deep Optimization kills services, resets settings after firmware updates. User must re-whitelist after every OTA. \\
Samsung   & 4/5 & Sleeping Apps list auto-adds apps not used in 3 days. Foreground services survive but battery optimisation aggressively throttles GPS. \\
Oppo      & 3/5 & ColorOS battery saver kills background apps. Less aggressive than MIUI but inconsistent across versions. \\
Vivo      & 3/5 & Similar to Oppo. Auto-start management restricts cold launch of services. \\
\bottomrule
\end{tabular}
\end{table}

The aggression scale is defined by three observable signals, each
testable at runtime via the Android \texttt{Build} class and
\texttt{PackageManager.resolveActivity()}:

\begin{itemize}[nosep]
  \item \textbf{S1: Foreground service kill.} The OEM kills a foreground
        service that holds \texttt{FOREGROUND\_SERVICE\_LOCATION} and is
        whitelisted via \texttt{isIgnoringBatteryOptimizations()}.
  \item \textbf{S2: Settings reset on OTA.} The OEM resets battery
        optimisation or autostart whitelist entries after a firmware
        update, requiring the user to re-whitelist.
  \item \textbf{S3: Background-only kill.} The OEM kills background
        services or defers their restart, but foreground services survive
        if whitelisted.
\end{itemize}

\noindent Rating: 5/5 = S1 + S2 (kills whitelisted foreground services
\emph{and} resets the whitelist on OTA). 4/5 = S1 only (kills
foreground services but whitelist persists across OTA). 3/5 = S3 only
(background services killed; foreground services survive if whitelisted).
2/5 = stock Android battery optimisation, no additional OEM
restrictions. 1/5 = AOSP behaviour, no vendor modifications.

The aggression ratings come from field deployment with 86 physical
devices across three deploying organisations (two large fleets of 45 and
35 devices, one smaller with 6). Xiaomi accounts for 41.9\% of the
physical fleet; combined with OPPO, Vivo, and Realme, 65\% runs on
OEMs rated aggression 3/5 or higher. The ratings are directional, not
statistically rigorous; the fleet is too small for confidence intervals
per OEM. But the session data tells its own story: of 159 tracked
sessions over two months, 65 (40.9\%) captured three or fewer waypoints
despite running for hours, and 25 (15.7\%) ran until server timeout
without a user-initiated end. The order-of-magnitude gap between
``working'' and ``silently dead'' sessions is consistent with
practitioner reports on dontkillmyapp.com, and the database cannot tell
us why because the telemetry to classify kill causes does not exist.

The only public resource for OEM-specific behaviour is
dontkillmyapp.com~\cite{dontkillmyapp}, which is maintained by the
developer of the Don't Kill My App benchmarking tool. It catalogues
workarounds per OEM but provides no detection methodology and no
severity data. Our methodology fills this gap: not by replacing the
workaround catalogue, but by giving apps a way to measure how often
each kill vector actually fires in their fleet.

\subsection{Permission Revocations}

Permission revocation is the quietest kill vector. On Android 12+, the
system auto-revokes permissions for apps it deems
inactive~\cite{android_hibernate}. Users also manually downgrade from
``Always'' to ``While using the app'' or revoke location entirely. In
both cases, there is no callback, no crash, no error thrown in the
running process. Location updates simply stop. The app continues running,
believing it is tracking, producing zero-waypoint gaps that are
indistinguishable from GPS hardware failure unless the app actively polls
its own permission state.

We observed this in production: a permission revocation during an active
tracking session produced no error. The app's distance counter froze at
the last known value. The tracked worker completed a 45 km route and
the management dashboard showed 0.0 km for the last three hours of the
session. The first indication of a problem was a phone call from the
deploying organisation, not an alert.

\subsection{Background Location Engine Architecture}
\label{sec:engine}

The audited app runs a background location engine compiled from Rust and
accessed from Flutter via FFI. The engine operates as an Android
foreground service with \texttt{FOREGROUND\_SERVICE\_LOCATION} (API 34+).
It maintains its own native odometer (cumulative distance in metres),
writes waypoints to local storage from a headless Dart isolate when the
UI process is not running, and registers a \texttt{BOOT\_COMPLETED}
receiver to restart after device reboot. The configuration sets
\texttt{stopOnTerminate: false} and \texttt{startOnBoot: true}.

This architecture produces three signals relevant to the measurement
methodology. First, a native odometer that resets to zero when the
service is fully killed and restarted: the reset is the fingerprint of a
process-level kill. Second, headless waypoints (marked
\texttt{source: 'headless'}) that prove the background service survived a
UI kill. Third, a \texttt{BOOT\_COMPLETED} event in the native layer
that fires on device restart. Of these three, only the headless waypoints
reach the backend. The odometer reset is logged to \texttt{debugPrint}
and the boot event produces no persisted record on the Dart side.

% ============================================================================
% 3. KILL-CAUSE TAXONOMY
% ============================================================================

\section{Kill-Cause Taxonomy}
\label{sec:taxonomy}

We classify background execution failures into five categories based on
their cause and the observable signals each leaves behind.
Table~\ref{tab:taxonomy} summarises the taxonomy.

\begin{table*}[t]
\caption{Kill-cause taxonomy: categories, observable signals, and detection requirements}
\label{tab:taxonomy}
\small
\begin{tabular}{@{}p{2.4cm}p{3.2cm}p{4.8cm}p{4.8cm}@{}}
\toprule
Category & Cause & Observable signals & Detection requirements \\
\midrule
\textsc{Doze Kill} &
Stock Android Doze / Light Doze suspends or kills the process &
Waypoint gap with duration matching Doze maintenance window cadence ($\sim$15 min intervals). Process restart with no user interaction. Battery level continuous (not charged during gap). &
Persisted \texttt{last\_clean\_exit} flag. Timestamp of last waypoint before gap. Device not rebooted (uptime continuous). \\

\textsc{OEM Kill} &
Vendor battery manager (PowerGenie, MIUI, etc.) terminates the foreground service or process &
Waypoint gap. Native odometer resets to zero (service restarted). No Doze maintenance pattern in gap timing. OEM identified as aggressive ($\geq$4/5 aggression rating). &
OEM and OS version attached to session. Native odometer checkpoint persisted to disk. Aggression rating lookup per OEM. \\

\textsc{User Force Stop} &
User swipes app from recents or taps ``Force Stop'' in Settings &
Process death. No headless waypoints during gap (headless isolate killed with app). Boot receiver does not fire (not a reboot). Gap starts at a non-Doze-aligned time. &
Headless task liveness signal (heartbeat). Persisted process-start marker with monotonic ID. \\

\textsc{Reboot} &
Device restarted (user-initiated, OTA, crash) &
Process death. Boot receiver fires (\texttt{BOOT\_COMPLETED}). Device uptime resets. Gap duration includes boot time. &
Boot receiver logs a \texttt{device\_reboot} event with timestamp. Previous uptime persisted. \\

\textsc{Permission Revocation} &
User or OS revokes location permission while app is running or between sessions &
Location updates stop but process may continue running. \texttt{checkPermission()} returns denied on next poll. No process death (unless combined with another kill). &
Periodic permission polling (not only on foreground resume). Permission state persisted to disk so that a revocation during process death is detectable on next launch. \\

\bottomrule
\end{tabular}
\end{table*}

Two notes on the taxonomy.

First, categories are not mutually exclusive. A Doze kill followed by
a permission revocation while the process is dead produces a compound
failure that must be decomposed. The classification procedure
(Section~\ref{sec:classification}) handles this.

Second, the taxonomy is conservative. We do not include low-memory kills
(\texttt{onTrimMemory} / OOM) as a separate category because we have not
observed them as a primary cause in our fleet. The devices are budget
handsets (Redmi 9A, 3 GB RAM) running a 104,886-line app. If
low-memory kills were frequent, we would expect to see them here. We
did not, which suggests that foreground service priority is sufficient
protection against OOM on current Android versions. This may not hold
for devices with 1--2 GB RAM or for apps without a foreground service.

% ============================================================================
% 4. INSTRUMENTATION SPECIFICATION
% ============================================================================

\section{Instrumentation Specification}
\label{sec:instrumentation}

This section defines the minimum telemetry a production app needs to
detect each failure category. The specification has three layers:
events, metadata, and persistence.

\subsection{Events}
\label{sec:events}

Seven events form the core of the measurement framework.
Table~\ref{tab:events} lists them with their trigger conditions and
required fields.

\begin{table}[h]
\caption{Core measurement events}
\label{tab:events}
\small
\begin{tabular}{@{}p{2.8cm}p{5.5cm}@{}}
\toprule
Event & Trigger and required fields \\
\midrule
\texttt{process\_start} &
App process created. Fields: \texttt{launch\_id} (monotonic),
\texttt{timestamp}, \texttt{was\_clean\_exit} (from persisted flag),
\texttt{device\_uptime}. \\

\texttt{process\_exit} &
\texttt{onDestroy} or clean shutdown path. Fields: \texttt{timestamp},
\texttt{exit\_reason}. Sets \texttt{last\_clean\_exit = true}
on disk. \\

\texttt{service\_restart} &
Foreground service re-created after being killed. Fields:
\texttt{timestamp}, \texttt{native\_odometer\_value} (0 indicates
full kill), \texttt{last\_known\_odometer}. \\

\texttt{permission\_change} &
Permission state differs from last persisted state. Fields:
\texttt{timestamp}, \texttt{previous\_state},
\texttt{current\_state}, \texttt{trigger} (poll / resume / boot). \\

\texttt{device\_reboot} &
\texttt{BOOT\_COMPLETED} receiver fires. Fields: \texttt{timestamp},
\texttt{previous\_uptime} (persisted), \texttt{current\_uptime}. \\

\texttt{heartbeat} &
Periodic liveness signal from background task. Fields:
\texttt{timestamp}, \texttt{battery\_level}, \texttt{gps\_state}. \\

\texttt{gap\_detected} &
Waypoint timestamp gap exceeds expected interval by $>$2x. Fields:
\texttt{gap\_start}, \texttt{gap\_end}, \texttt{gap\_seconds},
\texttt{classified\_cause} (filled by classification procedure). \\

\bottomrule
\end{tabular}
\end{table}

The \texttt{heartbeat} event deserves explanation. Its purpose is not
health monitoring in the distributed-systems sense. It is a liveness
proof: if heartbeats stop, the process is dead, and the last heartbeat
timestamp gives the time of death. Without it, the only evidence of a
kill is the absence of waypoints, which is ambiguous (the device could
be stationary, GPS could be off, the user could be indoors). A heartbeat
every 60 seconds adds negligible battery cost (no GPS, no network, just
a local timestamp write) and converts ``no data'' from ambiguous silence
into evidence of process death.

\subsection{Metadata}
\label{sec:metadata}

Every event must carry device context: OEM (\texttt{manufacturer}),
\texttt{model}, \texttt{os\_version}, \texttt{sdk\_int}. This context is
available from Android's \texttt{Build} class and costs nothing to
collect, but it must be attached at event time, not only at login. In
the production app we audited, device info was captured once at
authentication and never attached to tracking sessions or permission
events. An OS upgrade between login and failure would be invisible.

\texttt{process\_start} and \texttt{service\_restart} should also carry
the current battery optimisation state:
\texttt{isIgnoringBatteryOptimizations()}, the OEM-specific autostart
state (where detectable), and the app's current Standby Bucket (API
28+). These fields distinguish between kills where the app had proper
exemptions and kills where it did not, which is critical for separating
OEM aggression from misconfiguration.

\subsection{Persistence}
\label{sec:persistence}

The hardest requirement. The following state must survive process death:

\begin{itemize}[nosep]
  \item \texttt{last\_clean\_exit}: boolean, written on clean shutdown,
        cleared on dirty start.
  \item \texttt{last\_heartbeat\_timestamp}: last heartbeat time.
  \item \texttt{last\_permission\_state}: the permission state at last
        check.
  \item \texttt{native\_odometer\_checkpoint}: last known odometer value
        from the background service.
  \item \texttt{device\_uptime\_checkpoint}: uptime at last heartbeat.
  \item \texttt{process\_launch\_id}: monotonic counter incremented on
        each process start.
\end{itemize}

SharedPreferences (or equivalent key-value store) is sufficient. The
writes are small and infrequent. The critical constraint: these writes
must happen synchronously before the state they protect becomes relevant.
A heartbeat that writes the timestamp \emph{after} doing other work
defeats the purpose if the process is killed during that work.

One pattern we found in the audited app illustrates the failure mode.
The permission revocation state (\texttt{\_permissionRevokedAt},
\texttt{\_lastPermissionStatus}, \texttt{\_preRevokeSpeedMps}) was held
in static Dart fields, never persisted. If the OS killed the process
during a revocation gap, the gap was lost. On next launch, the
permission status field was null, so the first resume logged
\texttt{foreground\_resume} instead of \texttt{permission\_revoked}. The
revocation happened; the telemetry said it did not.

% ============================================================================
% 5. CLASSIFICATION PROCEDURE
% ============================================================================

\section{Classification Procedure}
\label{sec:classification}

Given the events from Section~\ref{sec:events}, the classification
procedure maps an observed gap to a kill cause. The procedure runs on
each detected gap after the app resumes or restarts.

\subsection{Decision Tree}

\begin{enumerate}
  \item \textbf{Check \texttt{device\_reboot} event.} If a reboot event
        exists with a timestamp inside the gap window: classify as
        \textsc{Reboot}. If the gap extends beyond boot completion,
        check for a secondary cause in the post-boot portion.

  \item \textbf{Check permission state.} If
        \texttt{last\_permission\_state} (persisted) was granted and
        current state is denied: classify as
        \textsc{Permission Revocation}. If the process also died (no
        heartbeats during gap), classify as compound:
        \textsc{Permission Revocation} + secondary kill cause.

  \item \textbf{Check \texttt{was\_clean\_exit}.} If the process
        restarted and \texttt{was\_clean\_exit} is false: the process
        was killed. Proceed to distinguish kill cause.

  \item \textbf{Check native odometer.} If the native background service
        odometer reset to zero: the service was fully killed and
        restarted, not just the UI process. This eliminates
        \textsc{User Force Stop} on devices where force-stop prevents
        boot-receiver restart (most OEMs). Classify as \textsc{OEM Kill}
        if the device OEM exhibits signal S1 (foreground service kill
        despite whitelist; aggression $\geq$4/5 per
        Section~\ref{sec:landscape}), or as \textsc{Doze Kill} if the
        gap duration matches Doze maintenance windows.

  \item \textbf{Check headless heartbeats.} If heartbeats continued
        during the gap: the background service was alive but the UI
        process was not. This is a UI-only kill (low memory or user
        swiped from recents without force-stopping the service).
        Classify as \textsc{User Force Stop} (partial).

  \item \textbf{No heartbeats, no odometer reset, no reboot.} The
        process died and the service did not restart. On OEMs exhibiting
        S1 or S3, classify as \textsc{OEM Kill}. On stock Android or Samsung,
        classify as \textsc{Doze Kill} if gap timing aligns with
        maintenance windows, otherwise \textsc{Unknown}.
\end{enumerate}

\subsection{Disambiguation Limits}

Some kill causes cannot be distinguished with client-side telemetry
alone. A Doze kill on a Xiaomi device could be Doze or MIUI; without
access to the OEM's battery manager logs, the classification is
best-effort. We accept this. The taxonomy does not promise forensic
certainty. It promises structured observation where none existed before.

The \textsc{Unknown} category is deliberate. A measurement system that
forces every observation into a named cause will misclassify more than
one that admits uncertainty. Our goal is a high-confidence subset (reboot
and permission revocation are unambiguous) plus a best-effort remainder.

% ============================================================================
% 6. VALIDATION
% ============================================================================

\section{Validation: Auditing a Production App}
\label{sec:validation}

We applied the methodology to Zontra, a production field workforce
tracking app built in Flutter. The app has 104,886 lines of Dart in
\texttt{lib/}, runs on 86 physical devices (94 registered, 8 emulators
excluded) across three deploying organisations (two primary fleets of 45
and 35 devices, one smaller with 6), and has been in
production for six months at the time of this audit. The database
analysis covers two months of tracked sessions (August 29 to October 29,
2026): 159 tracking sessions, 14,094 waypoints, 211 attendance sessions.
The app uses a Rust-based background location engine accessed via FFI
(Section~\ref{sec:engine}), with Firebase Authentication, Firestore for
real-time sync, and SharedPreferences for local persistence.

The fleet is 41.9\% Xiaomi (the most aggressive OEM at 5/5), with Vivo
(17.1\%), OPPO (10.5\%), and a long tail of 30+ distinct models from
ten OEMs. Android versions span 8.1 to 17. Two organisations account for
93\% of the fleet: one runs a controlled fleet of bulk-purchased POCO X8
Pro Max handsets, the other is BYOD across budget to mid-range devices.

We are also the developers of this app. This dual role is a threat to
validity (Section~\ref{sec:threats}) but an advantage for the audit: we
have full access to the source code, the production database, and the
incident history.

\subsection{Existing Detection Signals}

The app has six signals that partially detect background execution
failures:

\begin{enumerate}[nosep]
  \item \textbf{\texttt{PermissionEvent} model.} Fields:
        \texttt{timestamp}, \texttt{status}, \texttt{trigger},
        \texttt{previousStatus}, \texttt{gapDurationSeconds},
        \texttt{deadReckoningMeters}. Triggers:
        \texttt{tracking\_start}, \texttt{foreground\_resume},
        \texttt{permission\_revoked}, \texttt{permission\_restored},
        \texttt{app\_start}. Stored in SharedPreferences under
        \texttt{permission\_log}.

  \item \textbf{Native odometer regression.} If the native background
        service reports an odometer value lower than the last
        persisted checkpoint, the code infers the service was ``fully
        killed and restarted from 0.'' This detection exists in the
        codebase but is logged only to \texttt{debugPrint}. It does
        not produce a persisted event.

  \item \textbf{Waypoint gap inference.} Waypoints carry a
        \texttt{previousId} field. A break in the chain is documented
        in the model as ``permission revocation, process kill, or other
        interruption.'' The gap is observable but not classified.

  \item \textbf{Headless waypoints.} The background isolate writes
        waypoints with \texttt{source: 'headless'}. Their presence
        proves the background service was alive. Their absence during
        a gap is evidence (not proof) of process death.

  \item \textbf{Session summary absence.} The odometer writes a
        \texttt{OdometerSessionSummary} on clean teardown. A session ID
        with no summary is a killed session. This is never checked
        programmatically.

  \item \textbf{\texttt{tracking\_reconciled} event.} If the native
        service is tracking but the Dart engine's \texttt{\_running}
        flag is false on launch, a \texttt{tracking\_reconciled} event
        fires with \texttt{action=reattach}. This is the closest thing
        to a restart event and it is logged to Firebase Analytics.
\end{enumerate}

\subsection{Gaps Found}
\label{sec:gaps}

The audit found twelve gaps where the existing telemetry fails to detect
or classify failures. We group them by the instrumentation layer they
affect.

\textbf{Event gaps.}
(1)~No \texttt{process\_start} or \texttt{process\_exit} events exist.
The only launch signal is \texttt{app\_opened} with
\texttt{is\_cold\_start}, which does not carry a clean-exit flag or a
launch ID.
(2)~No \texttt{device\_reboot} event. The \texttt{BOOT\_COMPLETED}
receiver exists in the location engine's native code
(Section~\ref{sec:engine}) but produces no persisted event on the Dart
side.
(3)~No heartbeat. There is no periodic liveness signal from the
background service. Silence is ambiguous.
(4)~The native odometer regression (signal 2 above) goes to
\texttt{debugPrint} only. It is not persisted, not sent to analytics,
and not visible outside a debug build.

\textbf{Metadata gaps.}
(5)~Device context (OEM, OS version, SDK) is captured at login and never
attached to tracking sessions, permission events, or Sentry scope. An OS
upgrade after login is invisible.
(6)~Battery optimisation state is never logged. The event
\texttt{battery\_optimization\_exempt} is declared in the analytics
registry but has no call site.

\textbf{Persistence gaps.}
(7)~Permission revocation state is held in static Dart fields
(\texttt{\_permissionRevokedAt}, \texttt{\_lastPermissionStatus},
\texttt{\_preRevokeSpeedMps}) and lost on process death. A kill during a
revocation gap erases the revocation.
(8)~The gap duration is clamped to 60 seconds and dead-reckoning
distance to 5,000 metres before logging. Real gaps exceeding these
limits are truncated. A 3-hour gap from an OEM kill is recorded as 60
seconds.

\textbf{Detection gaps.}
(9)~Permission checks run only on foreground resume, triggered by
\texttt{didChangeAppLifecycleState} in a single widget
(\texttt{HomeBody}). There is no periodic poll, no background check. A
revocation during an active session is detected only when the user
returns to the home screen.
(10)~Only \texttt{locationWhenInUse} is checked. An ``Always'' to
``While using'' downgrade is invisible because both return granted for
\texttt{locationWhenInUse}.
(11)~The backend never receives permission events. The PRD specifies a
backend sink (documented in the API spec), but the app does not
upload the \texttt{permission\_log}. All permission data is local-only
and lost on app reinstall.
(12)~Silent \texttt{catch} blocks in lifecycle code
(\texttt{\_reattach}, \texttt{\_teardown},
\texttt{\_reconcileWithServer}) swallow errors. Failures in the
restart and recovery paths are not observable.

\subsection{Database Analysis}
\label{sec:database}

We analysed the production database to quantify what the gaps hide.
Table~\ref{tab:sessions} summarises the session-level findings.

\begin{table}[h]
\caption{Tracking session analysis (159 sessions, Aug--Oct 2026)}
\label{tab:sessions}
\small
\begin{tabular}{@{}lr@{}}
\toprule
Metric & Value \\
\midrule
Total tracking sessions & 159 \\
Total waypoints & 14,094 \\
Sessions with $\leq$3 waypoints & 65 (40.9\%) \\
Sessions ended by server timeout & 25 (15.7\%) \\
Sessions with location gap $>$30 min & 25 \\
Waypoint gaps $>$5 min (within sessions) & 509 \\
Longest intra-session gap & 110 h 50 min \\
Waypoint source: foreground & 81.7\% \\
Waypoint source: headless/background & 18.3\% \\
GPS off or unknown at session end & 100/159 (62.9\%) \\
Sessions ending $\leq$15\% battery & 5 (3.1\%) \\
\bottomrule
\end{tabular}
\end{table}

Three patterns stand out.

\textbf{Pattern 1: Sessions that start and immediately die.}
65 sessions captured three or fewer waypoints despite running for hours.
Session TRK-00001020 ran for 9 hours and 6 minutes and recorded one
waypoint. TRK-00001041 ran for 9 hours and 1 minute: also one waypoint.
These are sessions where the tracking engine started, captured an initial
fix, and then went silent. The existing telemetry cannot say why. Battery
was not the cause: TRK-00001044 ended at 96\% battery with two waypoints
over 9 hours.

\textbf{Pattern 2: Sequence resets as kill evidence.}
Three sessions show waypoint sequence numbers resetting to zero
mid-session: 182$\to$0 (with battery jumping from 47\% to 100\%,
indicating a kill during charging and a restart after), 174$\to$0, and
29$\to$0. These resets are the fingerprint of a full service kill and
restart. The app already detects this (signal 2, native odometer
regression) but logs it only to \texttt{debugPrint}. In the database,
these resets are visible only by inspecting raw waypoint sequences. One
session shows the inverse: sequence 39$\to$246, a gap of 207 missing
sequence numbers. These waypoints were generated by the headless
background service but never uploaded, accumulating locally until the app
reconnected.

\textbf{Pattern 3: The \texttt{permission\_event} table is empty.}
The backend schema defines a \texttt{permission\_event} table. It
contains zero rows. The client-side permission log is stored in
SharedPreferences and never uploaded. This confirms gap 11 from the
audit: all permission event data is local-only, lost on app reinstall,
and invisible to any server-side analysis.

\subsection{Incident Examples}

Two production incidents illustrate what the combined gaps cost.

\textbf{Incident 1: 45 km route recorded as 0.0 km.} A mid-session
permission revocation went undetected because the check runs only on
foreground resume. The tracked worker completed a route and the
management view showed zero distance for the last three hours. Time to
diagnose: 4 hours, because the permission log showed no revocation event
(gap 9).

\textbf{Incident 2: Session data missing on Xiaomi fleet.} MIUI killed
the foreground service on 8 of 12 Xiaomi devices during extended
sessions. The odometer regression detection fired (gap 4) but went to
\texttt{debugPrint}. There was no alert, no analytics event, no way to
distinguish OEM kills from other causes without SSH access to each
device. Time to diagnose: two days, most of it spent confirming that the
kills were OEM-caused and not a code regression.

Both incidents would have been caught within minutes if the
instrumentation specification from Section~\ref{sec:instrumentation} had
been implemented: heartbeat absence for Incident 1, persisted odometer
regression event with OEM metadata for Incident 2.

% ============================================================================
% 7. THREATS TO VALIDITY
% ============================================================================

\section{Threats to Validity}
\label{sec:threats}

\textbf{External validity.} The taxonomy and methodology are derived
from one production app on one platform (Android) with one use case
(field workforce tracking). Other app categories (health monitoring,
fitness, IoT) may encounter kill patterns we have not observed. The
OEM aggression ratings are based on 86 devices in one regional market;
OEM behaviour may differ across regions, firmware branches, and carrier
customisations.

\textbf{Internal validity.} The gap analysis combines a code audit with
a database analysis. We can count gaps (509 over 5 minutes) and identify
sessions that captured almost no data (65 of 159), but we cannot
classify the cause of each gap without the instrumentation this paper
proposes. The 40.9\% figure for near-empty sessions is a symptom count,
not a diagnosis. The incident examples are retrospective; we selected
them because they illustrate the gaps, which introduces selection bias.

\textbf{Construct validity.} The taxonomy classifies by cause. Some
practitioners may prefer classification by impact (data loss duration,
recovery success) or by user-facing symptom (distance gap, frozen
counter, session timeout). Our taxonomy is not the only valid one; it is
the one we found most actionable for debugging.

\textbf{Researcher bias.} We built the app, designed its telemetry,
and found it wanting. We have an incentive to frame the gaps as
significant. The mitigation: every gap is tied to a specific code
location and a specific detection failure that a reader can verify in
the described architecture.

\textbf{iOS.} We do not cover iOS. Apple's background execution model
is different (no foreground service equivalent, silent push throttling
at 2--3 per hour, 30-second background execution limit). The taxonomy
categories exist on iOS but the signals differ enough that a separate
treatment is warranted.

% ============================================================================
% 8. RELATED WORK
% ============================================================================

\section{Related Work}
\label{sec:related}

\textbf{Crash reporting and mobile observability.}
Firebase Crashlytics~\cite{firebase_crashlytics} and similar tools
(Bugsnag, Sentry) capture crashes, ANRs, and non-fatal exceptions.
They are the standard first layer of mobile reliability telemetry. But
they share a structural blind spot: a background process killed by the
OS produces no crash, no ANR, no exception. The kill is invisible to
the crash reporter because no code ran to report it. Our methodology
targets exactly this gap. Fan et al.~\cite{fan2018large} analysed
16,245 framework-specific exception traces across 2,486 Android apps
at ICSE~2018 and found that framework exceptions dominate crash reports.
Their work confirms that crash reporting is thorough for exceptions that
fire; it says nothing about silent kills where no exception fires at
all. Li et al.~\cite{li2020aging} conducted a large-scale study of
Android ANR and SNR events across 30,000 devices and found that
responsiveness problems are pervasive across hardware. ANR is the
closest the OS gets to reporting a background failure, but ANR requires
the process to be alive and unresponsive. A killed process is not
unresponsive; it is absent.

\textbf{Background execution, energy, and scheduling.}
Li et al.~\cite{li2014making} studied background energy drain on
Android. Pathak et al.~\cite{pathak2012energy} built eprof for
fine-grained energy profiling. Liu et al.~\cite{liu2016wake} studied
wake lock misuses across 44,736 Android apps and built a static
analysis tool to detect them. All three focus on energy cost, not on
whether the OS terminates the process. Gonz\'{a}lez-P\'{e}rez et
al.~\cite{gonzalez2022reliable} addressed the scheduling side directly:
their study of Android background task scheduling found that OS
restrictions make reliable periodic execution difficult, and they
released a scheduling library that works at up to one-minute intervals.
Their work is the closest to ours in motivation: they want background
tasks to run reliably. The difference is scope. They address scheduling
(when does the task run?); we address detection (did the task survive,
and if not, why?).

\textbf{Android permissions.}
Felt et al.~\cite{felt2012android} analysed the Android permission
model's effectiveness. Au et al.~\cite{au2012pscout} mapped the
permission specification. Wijesekera et
al.~\cite{wijesekera2015android} studied permission decisions in
context. This body of work treats permissions as a security and privacy
mechanism. We treat permission revocation as a reliability failure:
when a tracking app loses location permission mid-session, the security
model is working correctly, but the app's function is broken, and
detecting that break in real time is an unsolved measurement problem.

\textbf{Mobile app reliability.}
Liu et al.~\cite{liu2014characterizing} characterised performance bugs
in mobile apps. Nistor et al.~\cite{nistor2013discovering} studied
performance bugs that manifest only under specific usage patterns.
Neither addresses background execution kills as a failure category.
The gap is telling: the mobile reliability literature has studied
crashes, performance, and responsiveness, but silent process death is
not in the taxonomy of any prior study we found.

\textbf{OEM fragmentation and customisation.}
Wei et al.~\cite{wei2016taming} studied Android fragmentation at the
API level. The fragmentation that drives \textsc{OEM Kill} in our
taxonomy is different: it lives in vendor battery management policies,
not in API compatibility. Three studies address the OEM customisation
layer directly. Wu et al.~\cite{wu2015vendor} reversed vendor-customised
apps from 12 smartphones and found quality, reliability, and security
problems in nearly every one. Possemato et
al.~\cite{possemato2021trust} analysed 2,907 ROMs from 42 OEMs and
found that 20\% violated Android's Compatibility Definition Document.
Dong et al.~\cite{dong2024same} studied device-specific behaviours in
20,000+ apps and found that apps in Chinese third-party markets exhibit
more device-conditional code than Google Play apps, targeting
manufacturers with larger market share. None of these studies measures
OEM-caused process kills specifically. The only systematic
documentation of OEM battery management behaviour remains
dontkillmyapp.com~\cite{dontkillmyapp}, a practitioner resource.

\textbf{Flutter and cross-platform.}
Ghag~\cite{ghag2026testing} reported on testing a Flutter app with
heavy platform integration, documenting the coverage exclusion policies
and mock-boundary strategies required for the same codebase (Zontra).
That work addresses testing infrastructure; this work addresses runtime
observability. Bi{\o}rn-Hansen et al.~\cite{biorn2020flutter} surveyed
cross-platform approaches including Flutter but did not address
background execution or kill behaviour.

% ============================================================================
% 9. DISCUSSION
% ============================================================================

\section{Discussion}
\label{sec:discussion}

The methodology has a ceiling we have not broken through. Client-side
telemetry cannot distinguish between a Doze kill and an OEM kill on
devices where both could be responsible. The classification procedure
uses OEM aggression ratings as a heuristic, but a heuristic is not a
root cause. Breaking through would require access to the OS's battery
manager logs or a kernel-level probe, neither of which is available to a
production app.

We considered and rejected a controlled-experiment design as an
alternative validation strategy. Running each kill vector on a device
lab would produce clean causal data but would not validate the
methodology's ability to detect kills in production, where the
conditions are uncontrolled and the signals are noisy. The field
deployment, despite its small scale, tests exactly what matters: whether
the methodology detects real failures that real ad-hoc telemetry misses.

One finding surprised us. The most consequential gap was not a missing
event or a missing metadata field. It was a persistence failure:
in-memory permission state lost on process death (gap 7). This is a
pattern that would survive most code reviews. The variable names are
reasonable, the logic is correct when the process stays alive, and the
failure requires a specific interleaving (revocation, then kill, then
restart) that is hard to reason about in a code review but happens
routinely in the field. Measurement methodologies that focus on what
events to log will miss this class of bug. Persistence requirements are
the harder, less obvious part of the specification.

% ============================================================================
% 10. CONCLUSION
% ============================================================================

\section{Conclusion}
\label{sec:conclusion}

We have not solved the background kill problem. We have described how to
see it.

Seven events, a metadata contract, and a persistence specification
convert silent failures into classified observations. The audit of a
production app found twelve detection gaps; the database behind it held
509 unexplained waypoint gaps and 65 near-empty sessions out of 159. The
OEM Kill and Doze Kill boundary remains genuinely blurry on some devices,
and the taxonomy accepts that: five categories, one explicit
\textsc{Unknown}, no false confidence.

Two limitations bound these results. First, the methodology is validated
against one app on one platform in one use case. The taxonomy may be
incomplete for health monitors, fitness apps, or IoT controllers that
face different kill profiles. Second, we have not implemented the
specification and measured the before/after detection rate. The gap
analysis shows what existing telemetry misses; it does not yet show that
the proposed telemetry catches it. Both limitations point to the same
next step: implement the specification in a production app and measure
the change. We have not validated on iOS, where the background execution
model differs enough that the instrumentation does not transfer even
though the taxonomy categories do.

The measurement framework is methodology, not software. Any team with
a background-service app and a key-value store can implement it. The
harder part is not the code; it is recognising that the default
telemetry (crash reports, analytics events, permission checks on resume)
is structurally blind to the failures that cost the most.

% ============================================================================
% REFERENCES
% ============================================================================

\bibliographystyle{IEEEtran}
\bibliography{references}

\end{document}
